\documentclass[aps,prd,twocolumn,showpacs,superscriptaddress,groupedaddress,nofootinbib,english,floatfix]{revtex4-2}  

% Packages you will need
\usepackage{amsmath}  % Math formulas
\usepackage{amssymb}  % Math symbols (like \mathbb, \mathcal)
\usepackage{graphicx} % Images
\usepackage[pdfencoding=auto,hidelinks]{hyperref} % Hyperlinks
\usepackage{bm}       % Bold math
\usepackage{booktabs} % For professional table formatting
\usepackage{array}    % For better column definitions
\usepackage{tabularx, siunitx}
\usepackage[ngerman, main=english]{babel}
\usepackage{catchfilebetweentags} 
\usepackage{siunitx}
\sisetup{
    separate-uncertainty = true,  % Forces the ± symbol
    multi-part-units = single,    % Wraps units: (1.2 ± 0.1) kg
    sticky-per = true             % Helps with complex units
}
\usepackage{comment}
\usepackage[noabbrev]{cleveref} % Load this LAST in your preamble

\makeatletter
\let\l@de\l@ngerman
\makeatother
\makeatletter
\let\l@en\l@english
\makeatother

% Support compilation from either the repository root or paper/.
\IfFileExists{paper/0_abstract.tex}{
  \def\PMpaperdir{paper/}
  \def\PMcalculations{calculations/cosmology.tex}
}{
  \def\PMpaperdir{}
  \def\PMcalculations{../calculations/cosmology.tex}
}

\begin{document}

% Title Area
\title{Persistence Mechanics of the Quantum Vacuum VI: The Thermodynamics of Persistence at Cosmic Scale}


\author{Kate Lenore Meyer}
\affiliation{Independent Researcher}
\email{kate@meyerhome.net}

\date{\today}

\input{\PMpaperdir 0_abstract}

\maketitle % Generates the title block

% Main Content
\input{\PMpaperdir 1_introduction}

% Main Content
\input{\PMpaperdir 2_init}
\input{\PMpaperdir 3_tensions}
\input{\PMpaperdir 4_blackholes}
\input{\PMpaperdir 5_dark_sector}

\input{\PMpaperdir audit}
\input{\PMpaperdir conclusion}

\appendix
\section{Internal Consequences and Formal Verification}
The companion \texttt{lean} project links the new coordinate overhead directly to the earlier payload definition and verifies the revised rational outputs. It also proves exact interval propagation for fixed $\mathcal L$ and the invariant product linking the two late-time rescalings.

Two information results use finite substrate assumptions without an external cosmological field equation. First, there is no injective encoding of $C+1$ labels into $C$ slots, or of $C$ old states times more than one independent new input into the same $C$ slots. Retaining reversibility therefore requires additional storage or output outcomes. This does not specify a Hawking-radiation process. Second, independent equal-area cells have area-proportional log state count, with coefficient $\log(q)/a$; its physical normalization remains to be derived.

The additional modules establish reversible-label bounds, distinguish the two proposed CP inputs, derive local saturation remainders and convexity, check background dilution and continuity identities, and derive the comparison temperature from a supplied entropy law and first law. Paper VI also imports the earlier common-scale neutrino exclusion. These are conditional mathematical results; they do not supply the missing primordial state, cosmological metric, phase-space dynamics or radiation mechanism.

Running \texttt{./check.sh} in the Lean folder builds the library and audits every project declaration against the standard logical-axiom allowlist. These results distinguish internal consequences, explicit constitutive assumptions and open dynamical links. None of the latter is supplied by an unproved global physical axiom.

% Bibliography
% References are provided by the explicit bibliography below.

\begin{thebibliography}{9}
    \bibitem{Planck2018}
    Planck Collaboration, "Planck 2018 results. VI. Cosmological parameters," \textit{Astron. Astrophys.} 641, A6 (2020).
    
    \bibitem{Riess2022}
    A. G. Riess et al. (SH0ES Collaboration), "A Comprehensive Measurement of the Local Value of the Hubble Constant with 1 km/s/Mpc Uncertainty from the Hubble Space Telescope and the SH0ES Team," \textit{Astrophys. J. Lett.} 934, L7 (2022).
    
    \bibitem{DES-Y3}
    T. M. C. Abbott et al. (DES Collaboration), "Dark Energy Survey Year 3 Results: Cosmological Constraints from Galaxy Clustering and Weak Lensing," \textit{Phys. Rev. D} 105, 023520 (2022).
\end{thebibliography}


\end{document}
